\scanpage{36}
\unsure{On aura donc additivement}
\[
 H^*(Y)=\coprod_{\substack{i+j\leq n-1\\i,j\geq0}}L_X^jP^i(X)\quad\oplus\quad E^{n-1}\tag*{\unsure{$(\alpha)$|$(b)$}}
\]

L'opération
\[
\varphi_*:H^*(Y)\longrightarrow H^*(X)
\]
s'identifie \unsure{alors|donc} à $L_X$ sur le morceau $\coprod_{\substack{i+j\leq n-1\\i,j\geq0}}L_X^jP^{\unsure{j|i}}(X)$ \unsure{du premier morceau de|des premiers morceaux de} $(\alpha)$, (et est \deleted{\unsure{nul sur}} donc injectif sur ce morceau) et il est \uline{nul} sur le morceau $\deleted{\unsure{\coprod_{0\leq i\leq n-1} L_X^{n-1-i}P^i(X)\oplus|\coprod_{0\leq i\leq n-1}L_X^jP^i(X)\oplus}}\ E^{n-1}$, \deleted{qui correspond \unsure{par l'isom} \illegible{1} à $\sum_{n-1\leq i\leq2(n-1)}P^i(Y)$ \unsure{$*$}}. Ainsi on trouve le

\uline{Corollaire}
\[
\begin{aligned}
\operatorname{Ker}\varphi_* &\deleted{\unsure{=\sum_{n-1\leq j\leq2n-1}P^j(Y)=\operatorname{Ker}}}\ E^{n-1}(Y)\\
\operatorname{Im}\varphi_* &=\operatorname{Im} L_X=\bigcap_{0\leq i\leq n}\operatorname{Ker}p_i^{\unsure{X|\alpha}}
\end{aligned}
\]
(la dernière formule, car $H^*(Y)=\operatorname{Im}\varphi^*+E^{n-1}$ et $\varphi_*$ s'annule sur $E^{n-1}$, \unsure{donc} $\operatorname{Im}\varphi_*=\operatorname{Im}\varphi_*\varphi^*\deleted{\unsure{=I}}$, or $\varphi_*\varphi^*=L_X$).
\[
\left\{\begin{aligned}
p_i^X\varphi_*&=0&&\text{si }0\leq i\leq n\\
p_j^X\varphi_*&=\varphi_*p_{j-2}^Y&&\text{si }\unsure{n+1|\deleted{n+1}}\leq i\leq 2n.
\end{aligned}\right.
\]

\margin{\deleted{\unsure{fonctorialité|variation}}\\\unsure{variantes}}
On a
\[
\begin{aligned}
\left\{\begin{aligned}
\varphi^*L_X&=L_Y\varphi^*\\
L_X\varphi_*&=\varphi_*L_Y
\end{aligned}\right. \qquad&
\begin{aligned}
\varphi_*\varphi^*&=L_X,\quad\varphi^*\varphi_*=L_Y\\
&\text{\unsure{formule des projections|formule de projection|formules de projection}}
\end{aligned}
\end{aligned}
\]
\deleted{Pour la \unsure{deuxième|dernière} formule, on \unsure{note} que les deux \unsure{mbs} sont nuls sur $E^{n-1}$, donc il suffit de vérifier,
\[
 (L_X\varphi_*)\varphi^*=(\varphi_*L_Y)\varphi^*\quad\text{i.e.}\quad
 L_X(\varphi_*\varphi^*)=\varphi_*\underbrace{(L_Y\varphi^*)}_{\inserted{\varphi_*\varphi^*L_X}}=
\]
or les deux mbs sont $L_X^2$.}
% archive sheet number: 35

